%% =============================================================================
%% sm_f6_frequentist.tex — SM-E, Subsection E.6: Frequentist Contextualization
%% Body content only (\subsection level); parent structure in sm_e.tex
%% =============================================================================

\subsection{Frequentist contextualization}
\label{smf:frequentist}

A natural question is whether a simpler design-based analysis could
replace the full Bayesian hierarchical model.  We compare the HBB
pseudo-posterior estimates with conventional survey-weighted frequentist
estimates obtained from \texttt{survey::svyglm}
\citep{Lumley2004,Lumley2020} to (i)~validate the design-consistency
of the pseudo-posterior fixed effects and (ii)~clarify which model
features cannot be accommodated by the frequentist approach.

\paragraph{Frequentist specification.}
We fit separate margin models using the same five-covariate
specification as the HBB model.  For the extensive margin, we fit a
survey-weighted logistic regression using
\texttt{svyglm(\ldots, family = quasibinomial())} with $N = 6{,}785$
observations.  For the intensive margin, we fit a survey-weighted
linear regression of the IT enrollment share $y_i / n_i$ on the same
covariates, restricted to the $N^{+} = 4{,}392$ participants.  Both
models use the NSECE complex survey design (stratification, clustering,
and sampling weights) and produce linearization-based cluster-robust
standard errors.  No random effects, beta-binomial structure, or
cross-margin dependence can be incorporated.

\paragraph{Results.}
\Cref{smf:tab-frequentist} compares the survey-weighted svyglm
estimates with the M3b-W sandwich-corrected pseudo-posterior summaries.

\begin{table}[htbp]
\centering
\caption{Frequentist contextualization: survey-weighted svyglm vs.\
  Bayesian HBB (M3b-W, sandwich-corrected) point estimates and standard
  errors.  The extensive margin uses logistic link in both approaches;
  the intensive margin uses an identity link (svyglm) vs.\ logit link
  (HBB), so the scales differ.  SE Ratio $=$ HBB sandwich SE / svyglm SE.}
\label{smf:tab-frequentist}
% Generated by 86_B6_frequentist_comparison.R
\small
\begin{adjustbox}{max width=\textwidth}
\begin{tabular}{@{}l l rr rr r r@{}}
\toprule
 & & \multicolumn{2}{c}{svyglm (design-based)} &
   \multicolumn{2}{c}{HBB (sandwich-corrected)} & & \\
\cmidrule(lr){3-4}\cmidrule(lr){5-6}
Margin & Parameter & Est. & SE & Est. & SE & $\Delta$Est. & SE Ratio \\
\midrule
Extensive & Intercept & 0.662 & 0.063 & 0.764 & 0.044 & +0.103 & 0.70 \\
 & Poverty & -0.199 & 0.071 & -0.324 & 0.052 & -0.125 & 0.74 \\
 & Urban & 0.234 & 0.039 & 0.442 & 0.034 & +0.208 & 0.89 \\
 & Black & 0.086 & 0.075 & 0.478 & 0.044 & +0.392 & 0.59 \\
 & Hispanic & 0.059 & 0.074 & 0.006 & 0.051 & -0.054 & 0.70 \\
[3pt]
Intensive & Intercept & 0.483 & 0.005 & -0.242 & 0.015 & -0.725 & 3.02 \\
 & Poverty & 0.015 & 0.007 & 0.090 & 0.019 & +0.075 & 2.77 \\
 & Urban & -0.009 & 0.004 & -0.047 & 0.021 & -0.038 & 4.63 \\
 & Black & 0.034 & 0.008 & -0.017 & 0.018 & -0.050 & 2.30 \\
 & Hispanic & 0.010 & 0.006 & -0.097 & 0.017 & -0.107 & 2.92 \\
\bottomrule
\end{tabular}
\end{adjustbox}

\end{table}

For the extensive margin, both approaches use the logit link, enabling
a meaningful comparison on a common scale.  The coefficient signs and
relative magnitudes agree across all five predictors.  Point estimates
from HBB are systematically larger in absolute value than the svyglm
estimates, which is expected: the HBB model conditions on state-level
random effects, and conditional effects in mixed models are known to
exceed their marginal counterparts
\citep{ZegerLiangAlbert1988}.  The SE ratio
(HBB sandwich / svyglm) ranges from $0.59$ to $0.89$, reflecting the
variance reduction achieved by modeling between-state heterogeneity
through random effects rather than absorbing it into the residual.

For the intensive margin, the comparison is less direct because the two
approaches operate on different scales: svyglm regresses the
enrollment share $y_i / n_i$ under an identity link, while HBB models
the beta-binomial mean parameter $\mu_i$ under a logit link.  The
differing intercepts (svyglm: $+0.483$; HBB: $-0.242$) reflect this
scale difference rather than substantive disagreement, since
$\operatorname{logit}^{-1}(-0.242) \approx 0.44$, which is in the
range of the observed mean share.  The larger SE ratios for the
intensive margin ($2.30$--$4.63$) also arise from the logit-scale
parameterization, which inflates variances relative to the probability
scale.

\paragraph{Structural limitations.}
The comparison highlights several features of the HBB framework
that \texttt{svyglm} cannot accommodate:
\begin{enumerate}[nosep]
  \item \emph{Beta-binomial distribution:} svyglm models the continuous
    share $y/n$ via Gaussian errors, ignoring the bounded-count nature of
    the outcome and the extra-binomial overdispersion ($\kappa = 6.81$).
  \item \emph{Hurdle structure:} the two margins must be fitted
    separately, precluding joint inference on the zero-generating process.
  \item \emph{State-level random effects:} svyglm is a fixed-effect
    estimator; accommodating $S = 51$ state intercepts (and slopes in M3b)
    would require either fixed-effect dummies---inflating SE and
    sacrificing shrinkage---or a two-step approach that does not
    propagate uncertainty.
  \item \emph{Cross-margin correlation:} the bivariate random-effect
    structure $\boldsymbol{\delta}_s \sim \Norm_{2q}(\mathbf{0},
    \boldsymbol{\Sigma}_\delta)$ cannot be replicated in separate
    margin-specific regressions.
  \item \emph{Policy moderation:} the state-varying policy coefficients
    $\boldsymbol{\Gamma}$ interact hierarchical and design-based
    components in a way that has no natural frequentist counterpart.
\end{enumerate}

Taken together, the svyglm comparison serves as a fixed-effect-level
sanity check: the Bayesian pseudo-posterior fixed effects are
design-consistent with the frequentist target, while the full HBB
framework adds hierarchical structure, distributional flexibility,
and joint inference that the frequentist alternative cannot provide.
